\section{Discusión}

\subsection{Capacidad tecnológica conseguida}

El artefacto aporta la capacidad antes ausente: transforma una frase o un video con subtítulos en chunks trazables, compara tres enfoques, enlaza alertas con tiempos y registra correcciones. El clasificador Qwen3--LoRA, basado en un LLM, presenta la mayor estimación puntual de \(\AP\) y F1; E5 presenta el mayor recobrado de cualquier daño bajo los umbrales ordinarios; SVM ofrece una referencia más pequeña. La elección depende del costo de falsos negativos, revisión y cómputo. Además, las métricas agregadas no sustituyen las pruebas por fenómeno: trabajos como HateCheck y los análisis de sesgo de Sap \emph{et al.} muestran por qué conviene examinar conductas y grupos concretos \cite{rottger2021hatecheck,sap2019racialbias}.

La fusión operativa de acoso personal y amenaza directa aumenta el soporte y simplifica el contrato, pero pierde granularidad y no establece equivalencia conceptual o legal. Las etiquetas finas preservan esa distinción para auditoría y Qwen3--LoRA, aunque no deben presentarse como salidas operativas validadas. Del mismo modo, el consenso 2 de 3 es una regla implementada, no una mejora demostrada: necesita evaluación prospectiva.

\subsection{Alcance operativo conseguido y camino a producción}

El artefacto puede ensayarse fuera de línea o en modo sombra---sin actuar sobre usuarios---como moderador semiautomático: prioriza fragmentos, muestra el instante y el texto que originan cada alerta y entrega un diagnóstico preliminar a un supervisor. Este reparto entre modelo y persona es coherente con sistemas de moderación asistida y aprendizaje para diferir decisiones inciertas \cite{andersen2021rem,mozannar2020defer}. La política selectiva presenta en la Sección~VI una cobertura alta junto con su carga de revisión. Esa capacidad se consigue con modelos compactos, entrenamiento combinado entre CPU de consumo y Colab/CUDA, y un paquete de inferencia compatible con CPU; por ello el artefacto puede ensayarse sin poseer infraestructura GPU de centro de datos.

La derivación de seis de cada diez chunks cuantifica la capacidad humana que ese ensayo debe prever. Su recobrado corresponde al \emph{enrutamiento}, no al acierto en cada categoría; la Tabla~\ref{tab:final14} muestra el desempeño ordinario por separado. El test es enriquecido y sus rótulos son mayormente asistidos, por lo que todavía no estima prevalencia ni carga natural de YouTube; el cambio entre una muestra construida y la distribución de uso puede alterar el desempeño \cite{tonneau2024naijahate,moreno2012datasetshift}. El siguiente paso es un piloto prospectivo supervisado que mida latencia, costo, capacidad de revisión, seguridad y deriva en el hardware objetivo. El bloqueo, la sanción o la aprobación sin intervención humana no forman parte del propósito ni del alcance validado; la moderación que afecta usuarios requiere además reglas de gobierno y rendición de cuentas \cite{gorwa2020moderation}.
